\section{蒸馏、资源劣势与中国模型公司的机会}

谈到中美模型差距时，姚顺宇认为过去一年到一年半 gap 在缩小，但能否完全弥合甚至反超仍不确定。他特别提到，中国模型公司在算力资源上处于劣势，这种劣势可能逼出一些有趣策略，其中之一就是 distillation。

\begin{figure}[H]
\centering
\includegraphics[width=0.88\textwidth]{frame-distillation.jpg}
\caption{蒸馏章节：讨论“硬蒸”和更聪明的蒸馏方式。\protect\footnotemark}
\end{figure}
\footnotetext{视频画面时间区间：00:59:30--01:00:00。该帧对应硬蒸/软蒸的讨论场景。}

\begin{knowledgebox}{读图：蒸馏段落如何使用画面}
画面只能说明时间来源，不能说明蒸馏方法细节。蒸馏的技术信息需要读相邻表格：硬蒸、软蒸、环境蒸馏和自我蒸馏的区别在训练信号，而不是在视觉元素。
\end{knowledgebox}


\subsection{硬蒸与软蒸}

姚顺宇把蒸馏分成粗糙的“硬蒸”和更有科学含量的蒸馏。硬蒸就是直接拿大模型生成 token 强行训练，既有商业伦理问题，也显得智力上贫乏，因为它说明团队不知道自己真正想优化什么。更聪明的蒸馏则可能围绕能力迁移、偏好学习、数据生成、任务定义展开：不是抄答案，而是用强模型帮助构造更好的训练信号。

\noindent\begin{tabularx}{\textwidth}{p{0.18\textwidth}p{0.37\textwidth}X}
\toprule
类型 & 做法 & 风险/价值 \\
\midrule
硬蒸 & 直接收集强模型输出 token 做训练 & 法务和伦理风险高；容易学表面答案，缺少行为目标。 \\
软蒸 & 利用强模型辅助生成任务、反馈、解释或偏好信号 & 更接近科学问题：如何把能力、过程和评估迁移到弱模型。 \\
环境蒸馏 & 让模型在任务环境中产生轨迹，再筛选/评分 & 对 agent/coding 更有价值，但基础设施复杂。 \\
自我蒸馏 & 用模型自身或同族模型迭代提升数据 & 可能形成闭环，也可能放大错误和模式崩塌。 \\
\bottomrule
\end{tabularx}

\begin{importantbox}{蒸馏的本质不是复制答案，而是设计训练信号}
如果不知道自己要什么，蒸馏会退化成抄输出；如果清楚目标行为，蒸馏可以变成数据工程和能力迁移工具。两者差异在于：是否有明确任务定义、评估标准和失败分析。
\end{importantbox}

\subsection{机器人为什么仍有机会}

访谈中，姚顺宇对机器人实验室表达了比语言模型更强的兴趣，因为机器人仍有许多没有被做好的问题。语言模型已经不是蓝海；真正年轻的人如果想找机会，未必应该追逐最热方向，而应该找“没人做到好”的方向。机器人、多模态生成、量子调控等都可能是更蓝海的空间。

\begin{knowledgebox}{资源劣势可能改变研究品味}
当算力不能无限堆，团队会被迫重视效率、蒸馏、数据选择、训练信号和产品落地。资源劣势不是优势本身，但可能促使不同路径出现。中国模型公司的机会，可能正来自这种路径差异。
\end{knowledgebox}

\subsection{本章小结}

蒸馏不是一个简单的“偷不偷”问题，而是训练信号设计问题。硬蒸风险高且暴露目标缺失；软蒸和环境蒸馏则可能成为资源受限团队提升能力的重要路线。

